\documentclass[12pt,a4paper]{article}
\usepackage[utf8]{inputenc}
\usepackage[T1]{fontenc}
\usepackage[polish]{babel}
\usepackage{amsmath}
\usepackage{amssymb}
\usepackage{amsthm}

\begin{document}

\begin{center}
    \textbf{Kolokwium z Analizy Matematycznej 1.2* 2024/25} % [cite: 1] \\
    29 maja 2025, 14.30-17.30 % [cite: 3]
\end{center}

\vspace{0.5cm}

\noindent \textbf{Zadanie 1.} Oblicz granicę % [cite: 2]
\[
\lim_{n\to\infty}\sum_{k=n}^{2n}\frac{n^{2}}{n^{3}+kn^{2}+k^{2}n+k^{3}} % [cite: 4]
\]

\vspace{0.5cm}

\noindent \textbf{Zadanie 2.} Wykaż, że dla wszystkich $x\in\mathbb{R}$ zachodzi równość % [cite: 16]
\[
\int_{0}^{x}e^{-t^{2}/2}dt=\frac{f(x)}{g(x)} % [cite: 17]
\]
gdzie funkcje $f$ i $g$ definiujemy jako % [cite: 18]
\[
f(x):=x+\frac{x^{3}}{1\cdot3}+\frac{x^{5}}{1\cdot3\cdot5}+\dots, \quad \text{oraz} \quad g(x):=1+\frac{x^{2}}{2}+\frac{x^{4}}{2\cdot4}+\frac{x^{6}}{2\cdot4\cdot6}+\dots % [cite: 19]
\]

\vspace{0.5cm}

\noindent \textbf{Zadanie 3.} Rozważmy ciąg funkcyjny % [cite: 28]
\[
f_{n}(x):=n\cdot\left[e^{x}-\left(1+\frac{x}{n}\right)^{n}\right] \quad \text{dla } x\in[0,+\infty). % [cite: 29]
\]
Zbadaj zbieżność punktową, jednostajną i niemal jednostajną tego ciągu. % [cite: 30]

\vspace{0.5cm}

\noindent \textbf{Zadanie 4.} Rozważmy funkcję % [cite: 38]
\[
f(x)=\sum_{n=1}^{\infty}\frac{x}{n}\cdot e^{-x^{2}n}. % [cite: 39]
\]
Znajdź maksymalne podzbiory $\mathbb{R}$, na których funkcja $f(x)$ jest: % [cite: 40]
\begin{enumerate}
    \item[(a)] dobrze określona, % [cite: 41]
    \item[(b)] ciągła, % [cite: 42]
    \item[(c)] różniczkowalna, % [cite: 43]
    \item[(d)] klasy $C^{\infty}$. % [cite: 44]
\end{enumerate}

\vspace{0.5cm}

\noindent \textbf{Zadanie 5.} Rozważmy funkcję ciągłą $f:[a,b]\to\mathbb{R}$ i zbiór skończony $A\subset[a,b]$. % [cite: 66]
Udowodnij, że istnieje ciąg wielomianów $W_{n}$ jednostajnie zbieżny do funkcji $f$ na przedziale $[a,b]$ i taki, że $W_{n}(x)=f(x)$ dla każdego $x\in A$. % [cite: 67]

\end{document}